\documentclass{article}
\usepackage{amsmath,amssymb}
\usepackage{xurl}
\usepackage[hidelinks]{hyperref}
% Shared commands for notes in docs/paper-gaps/.

\DeclareUnicodeCharacter{2080}{\ifmmode{}_{0}\else\textsubscript{0}\fi}
\DeclareUnicodeCharacter{2081}{\ifmmode{}_{1}\else\textsubscript{1}\fi}
\DeclareUnicodeCharacter{2082}{\ifmmode{}_{2}\else\textsubscript{2}\fi}
\DeclareUnicodeCharacter{2099}{\ifmmode{}_{n}\else\textsubscript{n}\fi}

\newcommand{\Lean}{\textsc{Lean}}
\ExplSyntaxOn
\NewDocumentCommand{\leanid}{m}{
  \ifmmode
    \textsf{\detokenize{#1}}
  \else
    \group_begin:
    \urlstyle{sf}
    \tl_set:Nn \l_tmpa_tl {#1}
    \tl_replace_all:Nnn \l_tmpa_tl {\_} {_}
    \exp_args:NV \path \l_tmpa_tl
    \group_end:
  \fi
}
\ExplSyntaxOff
\providecommand{\tr}{\operatorname{tr}}
\newcommand{\tnleanIssueBase}{https://github.com/LionSR/TNLean/issues/}
\newcommand{\tnleanPrBase}{https://github.com/LionSR/TNLean/pull/}
\newcommand{\tnleanDocBase}{https://sirui-lu.com/TNLean/docs/}
\newcommand{\ghissue}[1]{\href{\tnleanIssueBase#1}{GitHub issue \##1}}
\newcommand{\ghpr}[1]{\href{\tnleanPrBase#1}{PR \##1}}
\newcommand{\doclink}[1]{\href{\tnleanDocBase#1.html}{\path{#1}}}

% Dirac notation and the identity operator, as in blueprint/src/macros/common.tex.
% \providecommand keeps note-local definitions authoritative.
\usepackage{dsfont}
\providecommand{\ket}[1]{|#1\rangle}
\providecommand{\bra}[1]{\langle #1|}
\providecommand{\braket}[2]{\langle #1 | #2 \rangle}
\providecommand{\ketbra}[2]{|#1\rangle\!\langle #2|}
\providecommand{\Id}{\mathds{1}}
\providecommand{\one}{\mathds{1}}
\providecommand{\C}{\mathbb{C}}
\providecommand{\MN}[1]{M_{#1}(\C)}
\providecommand{\MD}{M_D(\C)}

% Verdict marker: \gapnote{<kind>}{<status>}. Kind is the policy
% classification (clarification, local-correction, scope-restriction,
% unfaithful, false-source, open-gap); status is open, wip (elimination
% underway), resolved, or historical. Machine-read by the site index;
% prints a verdict line.
\newcommand{\gapnote}[2]{\par\noindent\textbf{Verdict:} #1 (#2).\par}

\title{Uniform Parent-Hamiltonian Gaps for Compact Normal Families}
\author{}
\date{2026-10-02}
\begin{document}
\maketitle
\gapnote{scope-restriction}{resolved}

\section{What the source states}
Appendix~A of \cite{Schuch2011Phases} proves that the family $H_\gamma$ of
parent Hamiltonians along the path from a matrix product state to its
isometric form has a spectral gap bounded below uniformly in
$\gamma\in[0,1]$ and in the chain length. The tensors of the path are in
normal form, $A_i(\gamma)=\bigoplus_\alpha A_i^\alpha(\gamma)$, possibly with
several blocks. The argument applies Nachtergaele's bound, controlled by the
second transfer eigenvalue of each block, by the overlaps of the local spaces
of different blocks, and by the gap $\Delta_{2p}(\gamma)$ of $H_\gamma$
restricted to $2p$ sites. The last quantity has a uniform lower bound
``as the restricted Hamiltonian is continuous in $\gamma$''.\footnote{Source
    traceability: \path{Papers/1010.3732/paper_v3.tex}, lines 2475--2580; the
    continuity step is lines 2575--2578.}

\section{The earlier restriction}
The first compact-family result applied only to one-site injective tensors
and the canonical two-site interaction. It did not cover normal forms with
several blocks, or the positive non-projector interactions of the source.
That restriction is removed by the results below. The proof uses finite-range
Knabe windows in place of the source's Nachtergaele estimate; this changes
the argument, not the uniform-gap conclusion. Compact families of normal
forms other than the source path are covered below only under a common
simultaneous spanning length or at the stated sufficient ranges; this is
not part of the source's statement.

\section{Compact normal-form families}
Let $x\mapsto(A_x^j)_j$ be continuous, with fixed positive block dimensions
$D_j$, and suppose the length-$S$ word tuples span the product matrix algebra
for every parameter, where $S>0$. Let $A_x$ be their block sum with nonzero
weights. The weights need not be continuous, since they do not affect any
local ground space. For every fixed $R\geq S+1$ and every compact parameter
set $K$, there is $\delta>0$ such that
\[
    \delta\lVert v\rVert\leq\lVert H_{N,R}(A_x)v\rVert
    \qquad(x\in K,\ N\geq R,\ v\perp\ker H_{N,R}(A_x)).
\]
This includes the minimal ring $N=R$. The joint boundary map identifies the
local ground space, while its restriction to scalar block boundaries
identifies the periodic kernel. Both maps are injective and continuous, so
their range projections are continuous. Compactness supplies uniform gaps
for finitely many short rings. Strict finite-range Knabe windows and a finite
subcover supply the common gap for all longer rings.

For pairwise inequivalent normalized normal blocks, a common simultaneous
spanning length is derived from the dimension bound. Any
$R\geq3\max(\sum_jD_j,1)^5$ suffices. For a single normal tensor, without a
normalization hypothesis, quantum Wielandt gives the sharper sufficient
range $D^4+1$. These bounds are sufficient, not asserted to be optimal.
\footnote{Formal declarations:
    \leanid{MPSTensor.exists_uniform_parentHamiltonianES_toTensorFromBlocks_gap_of_compact_all_lengths},
    \leanid{MPSTensor.exists_uniform_block_parentHamiltonianES_gap_of_compact_isNormalTensor}, and
    \leanid{MPSTensor.exists_uniform_parentHamiltonianES_gap_of_compact_isNormal_all_lengths}.}

\section{Positive interactions and the source path}
For a continuous family of positive local interactions $h_x$ with the same
local kernels, there is $\kappa>0$, uniform on a compact parameter set, such
that $h_x\geq\kappa P_{G_R(A_x)^\perp}$. Summing translated copies transfers
the canonical gap to these interactions. A pointwise upper comparison gives
equality of the periodic kernels; its constant need not be uniform.
\footnote{Formal declarations:
    \leanid{MPSTensor.exists_uniform_pos_parentInteractionES_le_of_compact},
    \leanid{MPSTensor.exists_uniform_periodicInteractionHamiltonianES_gap_of_compact_canonical_gap}, and
    \leanid{MPSTensor.exists_uniform_block_parentInteraction_gap_of_compact}.}

The source first blocks the tensor so that the joint one-site boundary map
is injective. Along its isometric deformation, this map is multiplied on
the physical space by the invertible map $Q_\gamma$. Simultaneous spanning
therefore persists with $S=1$. The preceding result applies at $R=2$ to the
source's continuous positive interactions. Thus the absence of several
blocks and non-projector interactions in the earlier result is no longer
a restriction.

The particular isometric deformation is also constructed from the original
block tensor. Let $P_{i,(j,a,b)}=A_i^j(a,b)$ and let $E$ be the orthogonal
projection onto $\operatorname{ran}P$. The left polar factors satisfy
$P=QW$, $Q=(PP^\dagger)^{1/2}+I-E>0$, and $WW^\dagger=E$.
The identity extension outside the retained physical space leaves the
factorization unchanged and keeps the original physical dimension.
Reading $W$ as block matrices gives a continuous path $Q_\gamma W$ whose
endpoints are $W$ and the original tensor. Its simultaneous word spans
are those of the original blocks, and both its canonical projections and
the source's positive interactions have a uniform periodic gap for all
$N\geq2$. For pairwise inequivalent normal blocks, the required initial
blocking is derived at a positive length $L$ satisfying
$L+1\leq3\max(\sum_jD_j,1)^5$.
\footnote{Formal declarations:
    \leanid{Matrix.leftPolarPos_mul_leftPolarIso},
    \leanid{Matrix.range_leftPolarIso_finalProjection},
    \leanid{MPSTensor.exists_uniform_isometricDeformation_parent_gap},
    \leanid{MPSTensor.exists_uniform_isometricDeformation_interaction_gap}, and
    \leanid{MPSTensor.exists_blocked_isometricDeformation_parent_gap_of_isNormalTensor}.}

For a one-site injective tensor, the deformation is also formalized through
the right polar factors. Writing its physical matrix as $M=VP$, the path
$M_\gamma=V((1-\gamma)I+\gamma P)$ is continuous and injective in the
original physical space, even when $M$ is rectangular. The positive factor
stays invertible throughout the interval. Unitary physical and virtual
covariance is preserved, and the canonical parent interactions retain the
on-site symmetry. This covers only the single-block case of
Section~II.C.\footnote{Formal declaration:
    \leanid{MPSTensor.exists_uniform_parentHamiltonianES_gap_polarDeformation}.}

The canonical parent projections along any continuous one-site injective
path with fixed on-site symmetry now also give a symmetric gapped
interaction path on the common physical space. The norm-gap estimate
gives a spectral gap above zero, since the periodic MPS vector is a
nonzero zero eigenvector. Applied to the polar path in unitary virtual
gauge, this joins the canonical parents of the original tensor and
its isometric form. It does not perform the different-dimension
endpoint embeddings in the classification theorem.\footnote{Affected
    declarations: \leanid{MPSTensor.canonicalInjectiveGappedPath} and
    \leanid{MPSTensor.polarGappedInteractionPath}.}

Normalization and unitary virtual covariance are now obtained from the
original one-site injective on-site symmetric tensor. Perron normalization
uses only nonzero rescaling and virtual gauge, which preserve all canonical
parent interactions. The positive definite adjoint fixed point is derived
from unitality by trace duality, and the virtual gauges are made unitary.
Thus the path to the isometric parent requires neither an assumed
normalization nor supplied virtual unitary matrices.\footnote{Affected
    declaration:
    \leanid{MPSTensor.exists_prepared_polarGappedInteractionPath_of_isOnSiteSymmetric}.}

The prepared path now also retains the virtual cohomology class of the
original tensor. A bond change of basis conjugates all virtual matrices
without changing their factor system. Unitary covariance choices for the
prepared tensor assemble into a single projective representation, which
implements symmetry at every point of the polar path. Gauge independence
then identifies the class with any endpoint choice.\footnote{Affected
    declaration:
    \leanid{MPSTensor.exists_prepared_polarGappedInteractionPath_with_virtual_class}.}

\section{General compact joint boundary inputs}
The joint boundary map is injective exactly when the homogeneous word tuples
span the product of the block matrix algebras. At a positive initial length,
this property persists at every larger length without normalization or
nonzero block-dimension assumptions. Openness and compactness therefore give
one positive threshold for a compact continuous family with pointwise
positive-length spanning. Trace preservation also allows an initial length
of zero by adjoining one letter. For normalized primitive blocks that are
pairwise inequivalent under similarity and phase, the pointwise hypothesis
follows without requiring continuous fixed-point matrices.

These statements supplement the compact normal-form gap conclusions above;
they do not replace those results by a weaker finite-window-only claim.
The weighted block family's local ground projection, canonical interaction,
and fixed-volume Hamiltonians are continuous at a common injectivity length.
Nonzero block weights need not vary continuously.
\footnote{Formal declarations:
    \leanid{MPSTensor.exists_uniform_wordTupleSpanTop_of_compact},
    \leanid{MPSTensor.wordTupleSpanTop_of_ge_of_pos},
    \leanid{MPSTensor.exists_uniform_blockGroundSpaceMapES_injective_of_compact_isPrimitiveMPS},
    \leanid{MPSTensor.continuous_iSup_groundSpaceES_starProjection_of_injective_family},
    \leanid{MPSTensor.continuous_openParentHamiltonianES_toTensorFromBlocks_family}, and
    \leanid{MPSTensor.continuous_parentHamiltonianES_toTensorFromBlocks_family}.}

\bibliographystyle{alpha}
\bibliography{references}
\end{document}
